\section{Scheduler and Running Example}
\subsection{Typed ownership and route resolution}
Each selected extension contributes immutable contract facets for fact ownership, pipeline effects, IR identities, and backend capabilities. Selection constructs one table for the actual runtime components. Duplicate fact owners, unknown facts, foreign production, impossible preservation, and ambiguous repeated module occurrences are rejected before semantic scheduling.

For every routable fact, the verifier registry resolves a triple $(rule,owner,b_{min})$. The rule names executable verification logic; the owner prevents an unrelated extension from satisfying the obligation with a similarly named check; and $b_{min}$ records the first boundary exposing the required artifact. Registry construction rejects conflicting owners and unknown boundaries. Effect validation rejects an invalidation that has no route, so a missing route cannot disappear before reaching the scheduler.

\subsection{Obligation creation and discharge}
When a selected transformation invalidates $f$, effect validation changes its state to invalid and emits $o=(f,rule,owner,b_c,b_e)$ with $b_e=\max(b_c,b_{min})$. The compilation observer retains the fact state and pending obligations while advancing monotonically through the boundaries. Obligations are normalized by deadline, rule, and fact. P2 selects exactly those whose first eligible boundary is current. It rejects an overdue obligation, an absent route, or an owner mismatch. P3 selects all boundary routes, but any due obligation remains attached to the corresponding invocation for lineage and diagnostics.

A successful production verifier invocation is the discharge event. It marks only facts owned by the invoked route and removes only matching due obligations; a baseline semantic check does not silently discharge a passive P1 obligation. The observer reports verifier diagnostics through the normal compiler error policy, and rejection prevents backend execution. Structural and semantic scopes are separate, avoiding a comparison in which a mandatory structural call silently performs every semantic check.

\subsection{Demand recomputation}
P1D receives an explicit set of queried facts. A query for an unknown fact fails closed. A query for an invalidated known fact schedules its canonical route. Without a query, P1D leaves the invalidation pending even if the pipeline reaches a semantic boundary. This is not an intentionally broken baseline: it models the normal guarantee of lazy analysis invalidation, where recomputation is demand-triggered rather than deadline-triggered.

\subsection{Deferred-route example}
The lifecycle is not limited to immediate post-pass checks. A production-path conformance case invalidates \texttt{backend.input-verified} at optimized AIR. Its canonical verifier requires the final AIR artifact and therefore declares backend input as its earliest executable boundary. P2 carries the pending obligation across the intervening observer return and discharges it in the newly exposed pre-backend hook. This case tests the same observer and scheduler used by normal compilation; it is an implementation-conformance witness, not an independently sampled compiler defect.

\subsection{Running example}
Figure~\ref{fig:running-example} follows one fault through the policies. A frontend produces AIR satisfying structural and selected capability facts. An optimizer performs a structurally legal rewrite to an intrinsic unsupported by the selected backend and invalidates \texttt{air.intrinsics-supported}. The generic AIR verifier still accepts the instruction schema.

P0 sees valid structure and proceeds. P1 records the invalidation but executes no semantic route. P1D rejects only in the variant where a downstream component explicitly queries the invalidated capability fact; in the matched unqueried variant it proceeds. P2 creates an obligation owned by \texttt{core.air} and executes it at optimized AIR, the first eligible boundary, before backend execution. P3 executes the same semantic verifier plus every other applicable route and yields the same correctness outcome on this case.

\begin{figure}[t]
\centering
\fbox{\begin{minipage}{0.94\columnwidth}
\small
source $\rightarrow$ frontend $\rightarrow$ valid AIR
$\rightarrow$ optimizer emits unsupported intrinsic
$\rightarrow$ optimized-AIR boundary
$\rightarrow$ backend\\[2pt]
\textbf{selected facts:} schema valid; branch targets valid; intrinsics supported\\
\textbf{effect:} invalidate intrinsics-supported\\
\textbf{P0/P1:} proceed \quad
\textbf{P1D:} reject iff queried \quad
\textbf{P2/P3:} reject at optimized AIR
\end{minipage}}
\caption{Composed pipeline and the selected contract relation in the running example. The rewrite is structurally legal; the violated relation is selected-backend capability support.}
\Description{A left-to-right compiler pipeline. An optimizer invalidates the selected backend capability fact. P0 and P1 proceed, P1D rejects only when queried, and P2 and P3 reject at the optimized-AIR boundary.}
\label{fig:running-example}
\end{figure}

\begin{figure}[t]
\centering
\fbox{\begin{minipage}{0.90\columnwidth}
\small
$\mathsf{valid}$ $\xrightarrow{invalidate}$ $\mathsf{invalid}$
$\xrightarrow{create\ o}$ $\mathsf{pending}$
$\xrightarrow[\text{sound accept}]{owner\ route}$ $\mathsf{discharged}$\\[-1pt]
\hfill $\searrow$ \textsf{rejected} (false fact, missing route, owner conflict, or missed boundary)
\end{minipage}}
\caption{Fact/obligation lifecycle. Meaning does not depend on color.}
\Description{A state transition diagram from valid fact to invalid fact, pending obligation, and either discharged or rejected.}
\label{fig:obligation-lifecycle}
\end{figure}
